\documentclass[10pt,a4paper]{amsart}
\usepackage[margin=19mm]{geometry}
\usepackage{amsmath,amssymb,mathtools}
\usepackage[scr=rsfs,cal=euler]{mathalfa}
\usepackage{xfrac}
\usepackage[unicode,hidelinks,bookmarks=false]{hyperref}
\newcommand{\ie}{i.e.\ }
\newcommand{\cf}{cf.\ }
\newcommand{\resp}{respectively\ }
\newcommand{\Subsection}{\subsection}
\providecommand{\nobreakdash}{\nobreak\hbox{-}}
\setlength{\emergencystretch}{2em}
\multlinegap0pt
\usepackage{amssymb}
\usepackage{mathtools}
\usepackage{xcolor}
\usepackage[shortlabels]{enumitem}
\newtheorem{theorem}{Theorem}
\newtheorem{lemma}{Lemma}
\newtheorem{proposition}{Proposition}
\newcommand{\Q}{\mathbb Q}
\newcommand{\Z}{\mathbb Z}
\title{On the sequence $\gcd(a^n-1,b^n-1)$}
\author{Khai-Hoan Nguyen-Dang}
\date{Source: arXiv:2606.07959v2 (CC BY 4.0)}
\begin{document}
\maketitle

\noindent\emph{Source and licence.} On the sequence $\gcd(a^n-1,b^n-1)$. Version arXiv:2606.07959v2 (CC BY 4.0), distributed by arXiv under the Creative Commons Attribution 4.0 International licence, http://creativecommons.org/licenses/by/4.0/, as recorded in the arXiv licence metadata for this exact version and shown on the abstract page https://arxiv.org/abs/2606.07959v2. Full licence notice: \texttt{licenses/arxiv-2606.07959-CC-BY-4.0.txt}.\par\smallskip

\noindent\textit{Editorial scope.} Counted unit: the article's theorem thm:common-factor (source0.tex lines 228--240). The article's complete original proof of it is delivered with it (source0.tex lines 570--638, one \texttt{proof} environment of 69 lines, which the article places away from the statement). No mathematics is written, repaired or re-derived: the statement and the proof are the article's own lines, in the article's own notation, and no retained line is substituted or re-flowed.  Retained uncounted as the source-local context the proof needs: lines 351--358; lines 363--370; lines 413--426; lines 473--490; lines 552--555.  Nothing was omitted from any retained span, so the package carries no omission record.  No retained line contains a citation, so the wrapper carries no bibliography and no bibliography entry.  No retained line contains a figure environment, an image-inclusion command or drawing code, so no image is delivered and the package declares no asset dependency.  Editorial formatting only, in addition: an AMS \texttt{amsart} preamble that restates the article's own declarations and macros used by the retained lines, namely the environments \texttt{theorem}, \texttt{lemma}, \texttt{proposition} on separate counters, together with the standard packages those lines need.  The retained environments print in this document as Theorem 1, Lemma 1, Proposition 1 in order of appearance, whereas the article prints the same items with the numbering of its own full document.  The complete native source of the article as distributed in the arXiv e-print for this exact version is bundled unchanged as \texttt{sources/202293/source0.tex}.
\begin{theorem}[Bugeaud--Corvaja--Zannier]\label{thm:BCZ}
Let $a,b\ge 2$ be multiplicatively independent integers.  Then, for every
$\varepsilon>0$, there is $n_0=n_0(a,b,\varepsilon)$ such that
\[
  \gcd(a^n-1,b^n-1)<\exp(\varepsilon n)
\]
for every $n\ge n_0$.
\end{theorem}

\begin{lemma}[Local valuation bound]\label{lem:valuation-bound}
Let $c>1$ be an integer and let $p$ be a prime.  There is a constant
$C_{c,p}$ such that
\[
  v_p(c^m-1)\le C_{c,p}+v_p(m)
  \qquad(m\ge 1).
\]
\end{lemma}

\begin{proposition}[Subexponential recurrences]\label{prop:quasipolynomial}
Let $(u_n)_{n\ge 1}\subset\Z$ be a sequence satisfying a constant-coefficient
linear recurrence, and assume that
\[
  \log^+|u_n|=o(n).
\]
Then there are an integer $M\ge 1$ and polynomials
$Q_0,\dots,Q_{M-1}\in\Q[X]$ such that, for every residue class
$r\pmod M$,
\[
  u_n=Q_r(n)
\]
for all sufficiently large $n\equiv r\pmod M$.
\end{proposition}

\begin{lemma}[Valuation-constrained quasipolynomials]\label{lem:local-polynomial}
Fix $N\ge 0$.  Let $(u_n)_{n\ge 1}$ be a sequence of nonzero integers.  Suppose
that there are $M\ge 1$ and $Q_0,\dots,Q_{M-1}\in\Q[X]$ such that
\[
  u_n=Q_r(n)
\]
for all sufficiently large $n\equiv r\pmod M$.  Suppose also that for every
prime $p$ there is a constant $C_p$ satisfying
\begin{equation}\label{eq:valuation-hypothesis}
  v_p(u_n)\le C_p+v_p(n+N)
  \qquad(n\ge 1).
\end{equation}
Then, for each $r\pmod M$, there are $\lambda_r\in\Q^\times$ and
$e_r\in\{0,1\}$ such that
\[
  Q_r(X)=\lambda_r(X+N)^{e_r}.
\]
\end{lemma}

\begin{lemma}\label{lem:eventual-periodic}
Let $(u_n)_{n\ge 1}$ satisfy a constant-coefficient recurrence with nonzero trailing
coefficient.  If $(u_n)$ is eventually periodic, then it is periodic.
\end{lemma}

\begin{theorem}\label{thm:common-factor}
Assume that $a$ and $b$ are multiplicatively independent, and fix $N\ge 0$.
Let $(W_n)_{n\ge 1}$ be an integer sequence satisfying a constant-coefficient
linear recurrence such that
\begin{equation}\label{eq:common-factor-hypothesis}
  W_n\mid a^{n+N}-1
  \qquad\text{and}\qquad
  W_n\mid b^{n+N}-1
  \qquad(n\ge 1).
\end{equation}
Then $(W_n)_{n\ge 1}$ is eventually periodic.  If $(W_n)$ satisfies a recurrence
with nonzero trailing coefficient, then it is periodic from the first term.
\end{theorem}

\begin{proof}[Proof of Theorem~\ref{thm:common-factor}]
Because $W_n$ divides the two nonzero integers in
\eqref{eq:common-factor-hypothesis}, one has $W_n\ne 0$ and
\[
  |W_n|\le \gcd(a^{n+N}-1,b^{n+N}-1).
\]
Applying Theorem~\ref{thm:BCZ} at the exponent $n+N$ gives
\[
  \log^+|W_n|=o(n+N)=o(n).
\]
By Proposition~\ref{prop:quasipolynomial}, there are $M\ge 1$ and polynomials
$Q_0,\dots,Q_{M-1}\in\Q[X]$ such that $W_n=Q_r(n)$ for all sufficiently large
$n\equiv r\pmod M$.

For every prime $p$, Lemma~\ref{lem:valuation-bound} gives
\[
  v_p(W_n)\le v_p(a^{n+N}-1)
  \le C_{a,p}+v_p(n+N).
\]
Lemma~\ref{lem:local-polynomial} now shows that
\begin{equation}\label{eq:linear-or-constant}
  Q_r(X)=\lambda_r(X+N)^{e_r},
  \qquad e_r\in\{0,1\},
\end{equation}
for every $r\pmod M$.

We claim that in fact $e_r=0$ for every $r$.  Suppose that $e_r=1$ for some
$r\in\{0,\dots,M-1\}$.  Put
\[
  d:=\gcd(r+N,M),\qquad r+N=d r_0,\qquad M=dM_0.
\]
Then $\gcd(r_0,M_0)=1$, except in the harmless case $r+N=0$, when $M_0=1$.
By Dirichlet's theorem, there are infinitely many primes
\[
  \ell\equiv r_0\pmod{M_0}.
\]
Write $\lambda_r=A_r/B_r$ in lowest terms, with
$A_r\in\Z\setminus\{0\}$ and $B_r\in\Z_{>0}$.  Choose such primes outside the
finite set dividing $a d A_rB_r$, and large enough that
\[
  n_\ell:=d\ell-N\ge 1
\]
lies in the eventual range.  Since $n_\ell\equiv r\pmod M$,
\eqref{eq:linear-or-constant} gives
\[
  W_{n_\ell}=\frac{A_r d\ell}{B_r}\in\Z.
\]
Because $\gcd(A_r,B_r)=1$ and $\ell\nmid B_r$, integrality forces
$B_r\mid d$.  Consequently
\[
  W_{n_\ell}=\left(\frac{A_r d}{B_r}\right)\ell,
\]
so $\ell\mid W_{n_\ell}$.  The divisibility hypothesis implies
\[
  \ell\mid a^{d\ell}-1.
\]
As $\ell\nmid a$, Fermat's theorem yields
$a^{d\ell}\equiv a^d\pmod\ell$, and hence
\[
  \ell\mid a^d-1.
\]
This is impossible for infinitely many distinct primes $\ell$.  Thus every
$e_r=0$.

Consequently, $(W_n)$ is eventually constant on each residue class modulo $M$,
so it is eventually periodic.  If $(W_n)$ satisfies a recurrence with nonzero
trailing coefficient, Lemma~\ref{lem:eventual-periodic} shows that it is periodic
from the beginning.
\end{proof}

\end{document}
